\documentclass[../../main.tex]{subfiles}

\begin{document}

\section{解的不存在性}\label{section:ch12-解的不存在性}

接下来识别半线性波动方程 $\square u+f(u)=0$ 不存在全局时间解的某些情形。方法是引入依赖于时间变量 $t$ 的适当积分量，然后导出涉及凸性的微分不等式，导致矛盾。

\subsection{负能量不存在性}\label{subsection:ch12-负能量不存在}

首先研究初值问题
\begin{align}
\begin{cases}
u_{tt}-\Delta u+f(u)=0 & \text{在 } \mathbb{R}^n\times(0,\infty) \text{ 中},\\
u=g,\;u_t=h & \text{在 } \mathbb{R}^n\times\{t=0\} \text{ 上}.
\end{cases}
\label{align:ch12-不存在初值}
\end{align}
总假设 $g,h\in C_c^\infty(\mathbb{R}^n)$ 且 $f(0)=0$。因此若光滑解存在，由 \S12.1.2\cref{theorem:ch12-依赖域}，它在每个时刻具有空间紧支集。

\begin{theorem}[负能量不存在性。]\label{theorem:ch12-负能量不存在}假设对某常数 $\lambda>2$ 有不等式
\begin{align}
zf(z)\leqslant\lambda F(z) \quad (z\in\mathbb{R}).
\label{align:ch12-zf不等式}
\end{align}
又设能量为负：
\begin{align}
E(0)=\int_{\mathbb{R}^n}\frac{1}{2}(|Dg|^2+h^2)+F(g)\,\mathrm{d}x<0.
\label{align:ch12-负初始能量}
\end{align}
则不存在对所有 $t\geqslant 0$ 存在的\eqref{align:ch12-不存在初值} 的光滑解 $u$。

因此 \S12.3.2\cref{theorem:ch12-三维短时间爆破} 证明中构造的解一般不能全局持续。
\end{theorem}

\begin{proof}
1. 定义
\begin{align*}
I(t)\coloneqq\frac{1}{2}\int_{\mathbb{R}^n}u^2\,\mathrm{d}x.
\end{align*}
则
\begin{align}
I''=\int_{\mathbb{R}^n}u_t^2+uu_{tt}\,\mathrm{d}x=\int_{\mathbb{R}^n}u_t^2+u(\Delta u-f(u))\,\mathrm{d}x=\int_{\mathbb{R}^n}u_t^2-|Du|^2-uf(u)\,\mathrm{d}x.
\label{align:ch12-I二阶}
\end{align}
分部积分有效，因 $u$ 对每个时刻具有空间紧支集。由能量守恒，
\begin{align*}
E(t)=\int_{\mathbb{R}^n}\frac{1}{2}(u_t^2+|Du|^2)+F(u)\,\mathrm{d}x=E(0)
\end{align*}
对每个 $t>0$。在\eqref{align:ch12-I二阶} 中加上减去 $(2+4\alpha)E(0)$：
\begin{align*}
I''=(2+2\alpha)\int_{\mathbb{R}^n}u_t^2\,\mathrm{d}x+2\alpha\int_{\mathbb{R}^n}|Du|^2\,\mathrm{d}x+\int_{\mathbb{R}^n}(2+4\alpha)F(u)-uf(u)\,\mathrm{d}x-(2+4\alpha)E(0).
\end{align*}

选取 $\alpha>0$ 使 $2+4\alpha=\lambda$，其中 $\lambda$ 是假设\eqref{align:ch12-zf不等式} 中的常数。则最后积分项非负，因此
\begin{align}
I''\geqslant(2+2\alpha)\int_{\mathbb{R}^n}u_t^2\,\mathrm{d}x-\lambda E(0).
\label{align:ch12-I二阶界}
\end{align}
因 $I'=\int uu_t\,\mathrm{d}x$，不等式\eqref{align:ch12-I二阶界} 蕴含
\begin{align}
(1+\alpha)(I')^2\leqslant(1+\alpha)\left(\int_{\mathbb{R}^n}u^2\,\mathrm{d}x\right)\left(\int_{\mathbb{R}^n}u_t^2\,\mathrm{d}x\right)\leqslant I(I''-\beta),
\label{align:ch12-Cauchy}
\end{align}
其中 $\beta\coloneqq-\lambda E(0)>0$。

2. 令
\begin{align*}
J\coloneqq I^{-\alpha}.
\end{align*}
由\eqref{align:ch12-Cauchy} 可计算
\begin{align}
J''=\alpha(\alpha+1)I^{-(\alpha+2)}(I')^2-\alpha I^{-(\alpha+1)}I''\leqslant-\alpha\beta I^{-(\alpha+1)}=-\alpha\beta J^{1+1/\alpha}.
\label{align:ch12-J二阶}
\end{align}
这表明 $J$ 是 $t$ 的凹函数。假设对某时刻 $t_0>0$ 有 $J'(t_0)<0$。则 $J$ 的凹性蕴含
\begin{align*}
J(t)\leqslant J(t_0)+(t-t_0)J'(t_0) \quad (t\geqslant 0),
\end{align*}
此不等式提供了大时刻 $t$ 时 $J(t)<0$ 的矛盾。

相反假设 $J'\geqslant 0$ 对所有 $t\geqslant 0$。则由\eqref{align:ch12-J二阶} 得
\begin{align*}
J''\leqslant-\alpha\beta J^{1+1/\alpha}(0)\coloneqq-\gamma.
\end{align*}
$\gamma>0$，因负能量假设\eqref{align:ch12-负初始能量} 蕴含 $g\not\equiv 0$。于是 $J'(t)\leqslant J'(0)-\gamma t<0$ 对大 $t$，再次矛盾。
\end{proof}

\subsection{小初值不存在性}\label{subsection:ch12-小初值不存在}

满足前节负能量假设\eqref{align:ch12-负初始能量} 要求 $g$ 不太小。值得注意的是，某些半线性波动方程即使对某些任意小且光滑的初值 $g,h$ 也不存在全局时间解。

作为例子，考虑三维空间中的初值问题：
\begin{align}
\begin{cases}
u_{tt}-\Delta u-|u|^p=0 & \text{在 } \mathbb{R}^3\times(0,\infty) \text{ 中},\\
u=g,\;u_t=h & \text{在 } \mathbb{R}^3\times\{t=0\} \text{ 上}.
\end{cases}
\label{align:ch12-负幂初值}
\end{align}
取任何光滑初值满足
\begin{align}
\int_{\mathbb{R}^3}g\,\mathrm{d}x>0,\quad\int_{\mathbb{R}^3}h\,\mathrm{d}x>0,
\label{align:ch12-初值正积分}
\end{align}
且 $\operatorname{spt}g,\operatorname{spt}h\subset B(0,R)$。我们将仅从这些条件推出不存在解，只要指数 $p>1$ 足够小。

\begin{theorem}[小数据不存在性。]\label{theorem:ch12-小数据不存在}假设
\begin{align*}
1<p<1+\sqrt{2}.
\end{align*}
则在上述 $g$ 和 $h$ 的条件下，初值问题\eqref{align:ch12-负幂初值} 不存在对所有 $t\geqslant 0$ 存在的光滑解 $u$。

此陈述应与 \S12.3 和 \S12.4 中对 $1\leqslant p\leqslant 5$ 证明的 $\square u+|u|^{p-1}u=0$ 在 $\mathbb{R}^3\times(0,\infty)$ 中光滑解的全局存在性形成对比。
\end{theorem}

\begin{proof}
1. 反设 $u$ 事实上是解并导出矛盾。因初值支集在 $B(0,R)$ 内，\S12.1.2\cref{theorem:ch12-依赖域} 蕴含 $u(\cdot,t)$ 支集在球 $B(0,R+t)$ 内。令
\begin{align*}
I(t)=\int_{\mathbb{R}^3}u\,\mathrm{d}x.
\end{align*}
则
\begin{align*}
I''=\int_{\mathbb{R}^3}u_{tt}\,\mathrm{d}x=\int_{\mathbb{R}^3}|u|^p\,\mathrm{d}x.
\end{align*}
因 $u(\cdot,t)$ 在球 $B(0,R+t)$ 外消失，
\begin{align*}
I=\int_{B(0,R+t)}u\,\mathrm{d}x\leqslant\left(\int_{\mathbb{R}^3}|u|^p\,\mathrm{d}x\right)^{1/p}|B(0,R+t)|^{1-1/p},
\end{align*}
因此
\begin{align}
I''\geqslant cI^p(1+t)^{-3(p-1)}
\label{align:ch12-Ip}
\end{align}
对常数 $c>0$。

2. 接下来引入线性波动方程的解 $v$：
\begin{align}
\begin{cases}
v_{tt}-\Delta v=0 & \text{在 } \mathbb{R}^3\times(0,\infty) \text{ 中},\\
v=g,\;v_t=h & \text{在 } \mathbb{R}^3\times\{t=0\} \text{ 上}.
\end{cases}
\label{align:ch12-辅助线性}
\end{align}
由推迟势公式\eqref{align:ch02-retarded}，
\begin{align*}
u(x,t)=v(x,t)+\frac{1}{4\pi}\int_{B(x,t)}\frac{|u|^p(y,t-|y-x|)}{|y-x|}\,\mathrm{d}y\geqslant v(x,t).
\end{align*}
现在\eqref{align:ch12-辅助线性} 和\eqref{align:ch12-初值正积分} 蕴含
\begin{align*}
\int_{\mathbb{R}^n}v\,\mathrm{d}x=c_1+c_2t
\end{align*}
对常数 $c_1,c_2>0$。此外，由于我们在 $n=3$ 空间维数工作，Huygens 原理告诉我们 $v$ 支集在环形区域 $A\coloneqq B(0,t+R)\setminus B(0,t-R)$ 内，其体积 $|A|\leqslant C(1+t)^2$。因此
\begin{align*}
c_1+c_2t=\int_A v\,\mathrm{d}x\leqslant\int_A u\,\mathrm{d}x\leqslant C\left(\int_{\mathbb{R}^3}|u|^p\,\mathrm{d}x\right)^{1/p}(1+t)^{2(1-1/p)}.
\end{align*}
于是
\begin{align*}
I''=\int_{\mathbb{R}^3}|u|^p\,\mathrm{d}x\geqslant c(1+t)^{2-p}
\end{align*}
对某常数 $c>0$。因 $I(0),I'(0)>0$，推出
\begin{align}
I\geqslant c(1+t)^{4-p}.
\label{align:ch12-I下界}
\end{align}

3. 设 $\varepsilon>0$，结合\eqref{align:ch12-Ip} 和\eqref{align:ch12-I下界}：
\begin{align*}
I''\geqslant cI^{1+\varepsilon}I^{p-1-\varepsilon}(1+t)^{-3(p-1)}\geqslant cI^{1+\varepsilon}(1+t)^{(4-p)(p-1-\varepsilon)}(1+t)^{-3(p-1)}=cI^{1+\varepsilon}(1+t)^{-\mu},
\end{align*}
其中 $\mu\coloneqq(p-1)^2+\varepsilon(4-p)$。因 $1<p<1+\sqrt{2}$，可固定 $\varepsilon>0$ 足够小使
\begin{align}
0<\mu<2.
\label{align:ch12-mu}
\end{align}
因 $I'>0$，
\begin{align*}
I''I'\geqslant cI'I^{1+\varepsilon}(1+t)^{-\mu},
\end{align*}
从而
\begin{align*}
((I')^2)'\geqslant c(I^{2+\varepsilon}(1+t)^{-\mu})'+c\mu I^{2+\varepsilon}(1+t)^{-\mu+1}\geqslant c(I^{2+\varepsilon}(1+t)^{-\mu})'.
\end{align*}
因此
\begin{align}
(I')^2(t)\geqslant(I')^2(0)+c(I^{2+\varepsilon}(t)(1+t)^{-\mu}-I^{2+\varepsilon}(0)).
\label{align:ch12-Ip2}
\end{align}
由\eqref{align:ch12-I下界}，对足够大的时刻 $t\geqslant t_0$，$I^{2+\varepsilon}(t)(1+t)^{-\mu}\geqslant 2I^{2+\varepsilon}(0)$。因此由\eqref{align:ch12-Ip2} 推出
\begin{align*}
(I')^2\geqslant cI^{2+\varepsilon}(1+t)^{-\mu}
\end{align*}
只要 $t\geqslant t_0$。然后
\begin{align*}
(I^{-\varepsilon/2})'=-\frac{\varepsilon}{2}I^{-\varepsilon/2-1}I'\leqslant-\frac{c\varepsilon}{2}(1+t)^{-\mu/2}.
\end{align*}
积分得
\begin{align*}
0\leqslant I^{-\varepsilon/2}(t)\leqslant I^{-\varepsilon/2}(t_0)-\frac{c\varepsilon}{2}\int_{t_0}^t\frac{\mathrm{d}s}{(1+s)^{\mu/2}}.
\end{align*}
这是矛盾，因由\eqref{align:ch12-mu} $\mu<2$，右端积分当 $t\to\infty$ 时发散。
\end{proof}

\end{document}